\begin{frame}{Stochastic simulations}
    \begin{wideitemize}
    \item Stochastic (with shocks) $Ef\left(x_{t-1},x_t,x_{t+1},u_t | \theta \right) =0,$ with decision rules $x_t = g(x_{t-1},u_t|\theta)$ and assumptions about shock distributions.
    \item Here, we search for decision functions: $g\left(\cdot\right)$ is unknown and can only be solved analytically in a few very simple cases.
    \item Dynare applies a \hi{perturbation} approach to solve model: A Taylor expansion of the original model can be used to obtain a Taylor expansion of the solution.
    \item To obtain the decision rules or policy function:
    \begin{equation*}
        x_t - x_0 = A\left(x_{t-1}-x_0\right)+Bu_t
    \end{equation*}
    \end{wideitemize}
\end{frame}

\begin{frame}[fragile]{Linearisation}
    \begin{wideitemize}
    \item We have linearised our three-equation model by hand.
    \begin{itemize}
        \item Declare the model as linear.
        \begin{minted}{c++}
        model(linear);
        \end{minted}
    \end{itemize}
    \item Dynare can linearise the model for you! Use it if you can.
    \item The model needs to be stationary to allow for linearisation around the steady state
     \begin{itemize}
        \item Dynare can handle deterministic and stochastic trends.
        \item One can always add trends back to the model.
    \end{itemize}
    \item Always (try to) provide an analytical steady state. However, Dynare can numerically solve for it based on some starting value (\texttt{initval;}).
    \end{wideitemize}
\end{frame}


\begin{frame}{Our model}
\begin{align*}
y_t 	=& E_t \left[y_{t+1} \right]- \frac{1}{\sigma}\left(i_t-E_t\left[\pi_{t+1}\right]\right)+\varepsilon_t^c \\
\pi_t 	=& \beta E_t\left[\pi_{t+1}\right]+\kappa y_t +\varepsilon_t^a \\
i_t 	=& \phi_\pi \pi_t + \phi_y y_t + \varepsilon_t^i
\end{align*}
in matrix form
\begin{equation*}
\underbrace{
\begin{bmatrix}
-1 & -1 & 0\\
-0 & \beta & 0 \\
0 & 0 & 0
\end{bmatrix}}_{A_1}
\underbrace{
\begin{bmatrix}
E_t\left[y_{t+1} \right]\\
E_t\left[\pi_{t+1} \right]\\
E_t\left[i_{t+1} \right]
\end{bmatrix}}
_{E_t[X_{t+1}]}
=
\underbrace{
\begin{bmatrix}
-1 & 0 & -\frac{1}{\sigma}\\
\kappa & -1 & 0 \\
\phi_y&\phi_\pi & -1
\end{bmatrix}}_{A_0}
\underbrace{
\begin{bmatrix}
y_{t}\\
\pi_{t}\\
i_{t}
\end{bmatrix}}_{X_t}
+\underbrace{
\begin{bmatrix}
1 & 0 & 0\\
0 & 1 & 0\\
0 & 0 & 1
\end{bmatrix}}_{\mathcal{E}}
\underbrace{
\begin{bmatrix}
\varepsilon^c_{t}\\
\varepsilon^a_{t}\\
\varepsilon^i_{t}
\end{bmatrix}}_{u_t}
\end{equation*}
\begin{itemize}
    \item Note: We used a slightly different notation by including the shock processes in matrix $\mathcal{E}$. In Dynare, $u_t$ contains only the white noise for stochastic variables, i.e. our $u$'s.
\end{itemize}
\end{frame}

\begin{frame}{How to solve the model (linearly)}
\begin{wideitemize}
\item We have a system of the form
\begin{equation*}
A_1E_t\left[X_{t+1}\right] + A_0X_t  + A_{-1}X_{t-1} +  \mathcal{E} u_t = 0
\end{equation*}
where in our case $A_{-1}$ is a matrix of zeros.
\item There are many algorithms to solve this system.
\item Dynare uses a generalised Schur decomposition based on the method of Klein (2000)\footcite{klein2000using} to obtain the solution of the form:
\begin{equation*}
    X_{t} = TX_{t-1}+Cu_t
\end{equation*}
\item Villemot (2011) shows you the algebra.\footcite{villemot2011solving}
\end{wideitemize}
\end{frame}


\begin{frame}[fragile]{Stochastic simulation command}
    \begin{wideitemize}
    \item The command  Simulates data from the model and performs IRFs \href{https://www.dynare.org/manual/the-model-file.html?highlight=stoch_sim#stoch_simul}{\texttt{stoch\_sim}}
    computes a Taylor approximation around the steady state.
    \item \texttt{order} sets the approximation order.
    \item \texttt{period} sets the number of simulation periods
    \item You can select variables to be plotted and printed on the screen.
    \begin{minted}{c++}
    stoch_simul(irf=20, order=1,periods=100) y;
    \end{minted}
    \item The Dynare reference manual lays out the many available options.
    \end{wideitemize}
\end{frame}

\begin{frame}{Accessing results}
\begin{wideitemize}
\item You can access the results of the solution (decision rules) in the results structure \texttt{oo\_.dr}
\item Field \texttt{ys} contains the steady state.
\item Field \texttt{ghx} contains the transition matrix (with all endogenous variables in deviation from the steady state)
\item Field \texttt{ghu} contains the matrix attached to the exogenous shocks
\end{wideitemize}
\end{frame}


\begin{frame}{Hands-on}
\begin{wideitemize}
\item Let's look at the \hi{code for stochastic simulations} in practice.
\end{wideitemize}
\end{frame} 